

%\illustration{0.3}{windmill.jpg}{
    %One of the flour mills.
    %CC BY-ND 2.0 by Ryk Venema on \href{https://www.flickr.com/photos/89113252@N00/240993655/}{Flickr}
%}
\illustration{0.3}{IMG_4445_small.jpg}{
    One of the flour mills.
}
% Photo taken by Mike de Vries




It is the 25th of September, 1825.
Having just completed a record-setting wheat harvest,
you wonder what to do with all this wheat.
With a stroke of genius, you decide to use all of this wheat to bake your favourite type of cookie: stroopwafels.
Of course, all this wheat should be ground to flour first.
Because you cannot wait to start baking, you want to do this as fast as possible,
so you decide to contact all flour mills in the Netherlands to ask them for help.

For each of these flour mills,
you know how fast they can grind wheat to flour
and how long it takes for a shipment to travel to and from the mill.
You have enough grain ships available that can transport wheat to these mills,
and bring the flour back from the mills.
Dividing the wheat optimally between these mills,
how long does it take until you have received back all the wheat?

As an example, consider the first sample case.
To divide the wheat optimally between the three mills,
you ship $400$ kilograms to the first, $120$ to the second, and $480$ to the third.
The first mill requires $5$ hours to grind its wheat,
the second requires $1$ hour, and the third requires $3$ hours.
Combined with the shipping time to and from each mill,
you receive all wheat back after exactly $11$ hours.

In the second sample case, we send all of the wheat to the first mill.
This mill can grind all $100$ kilograms of wheat in $1$ hour,
which together with the $2$ hours for shipping back and forth results in a total of $3$ hours.
As the shipping for the second mill would already take $4$ hours, it is optimal to only use the first mill.

